% arara: clean: {
% arara: --> extensions:
% arara: --> ['aux', 'log', 'out', 'pdf']
% arara: --> }
% arara: lualatex: {
% arara: --> shell: no,
% arara: --> draft: yes,
% arara: --> interaction: batchmode
% arara: --> }
% arara: lualatex: {
% arara: --> shell: no,
% arara: --> draft: no,
% arara: --> interaction: batchmode
% arara: --> }
% arara: clean: {
% arara: --> extensions:
% arara: --> ['aux', 'log', 'out']
% arara: --> }
\documentclass{scrletter}
\usepackage{graphicx}
\usepackage[shortlabels]{enumitem}
\usepackage{mathtools}
% \usepackage{plantuml}
\usepackage{hyperref}

\usepackage[
	citestyle=numeric,
	style=numeric,
	backend=biber,
]{biblatex}

\addbibresource{references.bib}

\newcommand{\MVAt}{{\usefont{U}{mvs}{m}{n}\symbol{`@}}}
\renewcommand{\thesection}{\Roman{section}}
\renewcommand{\thesubsection}{\thesection.\Roman{subsection}}


\begin{document}
\begin{center}
    \sffamily\bfseries\Large
    Estadística Inferencial~CM 2H2 A\\
    Prueba de entrada
\end{center}

\begin{center}
    \fbox{\fbox{\parbox{5.5in}
            {\centering
                Responda las preguntas en los espacios provistos en
                las hojas de preguntas.
                Los argumentos y la claridad de las respuestas se
                considerarán en la puntuación final.
            }}}
\end{center}
\vspace{0.1in}
\makebox[\textwidth]{Nombres y apellidos del alumno:\enspace\hrulefill}
\vspace{0.2in}
\providecommand{\name}{Fidel Jara Huanca}
\makebox[\textwidth]{Nombres y apellidos del profesor:\enspace \name.\hfill}

\begin{questions}
    % \header{Math 115}{Second Exam}{July 4, 1776}
    \question

Dada la siguiente tabla
\begin{table}[ht!]
    \centering
    \begin{tabular}{|C|C|C|C|C|C|}
        \hline
        x_{i}               & 2            & 5           & 7           & 11          & 20           \\
        \hline
        f\left(x_{i}\right) & \frac{1}{16} & \frac{1}{8} & \frac{5}{8} & \frac{1}{8} & \frac{1}{16} \\
        \hline
    \end{tabular}
\end{table}

\begin{parts}
    \part[5]

    Halle el valor de
    \begin{math}\displaystyle
        \mu=
        \sum_{i=1}^{5}
        x_{i}f\left(x_{i}\right)
    \end{math}.

    \part[5]

    Analice la convergencia de
    \begin{math}\displaystyle
        I=
        \int_{0}^{\infty}
        xe^{-\big(\frac{x-\mu}{2}\big)^{2}}\dl x.
    \end{math}
\end{parts}

% \begin{figure}[ht!]
%     \centering
%     \begin{tikzpicture}[xscale=0.5,yscale=0.5]
%         \draw[help lines] (0,0) grid (32,28);
%     \end{tikzpicture}
% \end{figure}

\begin{solutionordottedlines}
    \begin{parts}
        \part

        \textbf{Cálculo de} $\mu$.

        \begin{math}
            \mu=
            \sum_{i=1}^{5}
            x_{i}f\left(x_{i}\right)=
            2\times\dfrac{1}{16}+
            5\times\dfrac{1}{8}+
            7\times\dfrac{5}{8}+
            11\times\dfrac{1}{8}+
            20\times\dfrac{1}{16}=
            \dfrac{124}{16}=
            7.75
        \end{math}.

        \part

        \textbf{Convergencia de la integral}.

        Si la sustitución es $z=\frac{x-\mu}{2}$, entonces
        $x=2z+\mu$ y $\dl z=\frac{\dl x}{2}$.

        \begin{equation*}
            I=
            \int_{0}^{\infty}
            xe^{-\big(\frac{x-\mu}{2}\big)^{2}}
            \dl x=
            \int_{0}^{\infty}
            \left(2z+\mu\right)e^{-z^{2}}2\dl z=
            4\int_{0}^{\infty}
            ze^{-z^{2}}\dl z+
            2\mu\int_{0}^{\infty}e^{-z^{2}}\dl z<
            \infty.
        \end{equation*}

        Como cada integral converge, $I$ es convergente.
    \end{parts}
\end{solutionordottedlines}
    \question

Sea $X$ una variable aleatoria continua con función de densidad
\begin{math}
	f_{X}\left(t\right)=
	\begin{cases}
		0,                 & \text{si }t\leq 0.   \\
		k\left(1+t\right), & \text{si }0<t\leq 1. \\
		\frac{2}{3},       & \text{si }1<t\leq 2. \\
		0,                 & \text{si }t>2.
	\end{cases}
\end{math}

\begin{parts}
	\part[2]

	Determine el valor de $k>0$ para que $f_{X}$ sea una función de
	densidad.

	\part[3]

	Calcule $\mathbb{P}\left[\frac{1}{2}<X\leq\frac{3}{2}\right]$.
\end{parts}

\begin{figure}[ht!]
	\centering
	\begin{tikzpicture}[xscale=0.5,yscale=0.5]
		\draw[help lines] (0,0) grid (32,16);
	\end{tikzpicture}
\end{figure}

    \question

Sean $\Omega$ el espacio muestral y $A,B\subset\Omega$ dos sucesos
tales que $\mathbb{P}\left(A\right)=\frac{1}{4}$,
$\mathbb{P}\left(B\mid A\right)=\frac{1}{2}$ y
$\mathbb{P}\left(A\mid B\right)=\frac{1}{4}$.

Indique y justifique si las siguientes afirmaciones son verdaderas o falsas.

\begin{parts}
	\part[2]

	$A$ y $B$ son mutuamente excluyentes.

	\part[2]

	$\mathbb{P}\left(\overline{A}\mid\overline{B}\right)=\frac{3}{4}$.

	\part[2]

	$\mathbb{P}\left(A\mid B\right)+\mathbb{P}\left(A\mid\overline{B}\right)=1$.
\end{parts}

\begin{figure}[ht!]
	\centering
	\begin{tikzpicture}[xscale=0.5,yscale=0.5]
		\draw[help lines] (0,0) grid (32,12);
	\end{tikzpicture}
\end{figure}

		\question

Una máquina produce piezas del tipo $T$
para una clase de mueble, la probabilidad de que cualquier pieza
del tipo $T$ producida por la máquina resulte de buena calidad es $0.95$,
si se escoge una muestra aleatoria de $10$ de dichas piezas
para someterla a control de calidad.
Suponga que la calidad de cada pieza es independiente de las demás.

\begin{parts}
	\part[1]

	¿Cuál es la probabilidad de que $8$ resulten de buena calidad?

	\part[1]

	¿Cuál es la probabilidad de que todas resulten de buena calidad?

	\part[1]

	¿Cuál es la probabilidad de que por lo menos $8$ resulten de buena calidad?

	\part[1]

	¿Cuál es la probabilidad de que a lo sumo $7$ resulten de buena calidad?
\end{parts}

\begin{figure}[ht!]
	\centering
	\begin{tikzpicture}[xscale=0.5,yscale=0.5]
		\draw[help lines] (0,0) grid (32,16);
	\end{tikzpicture}
\end{figure}

\end{questions}
\vfill
\begin{flushright}\bfseries
    Facultad de Ciencias\\[2mm]
    Universidad Nacional de Ingeniería\\[2mm]
    28 de agosto de 2025%\today %\unmarkedfntext{Cada pregunta vale 5 puntos.}
\end{flushright}
\end{document}
